% GENERATED FILE. Edit canonical lessons, then run npm run generate:books.
% canonical-source-hash: fnv1a64:7082026a990aa35d
% canonical-lessons: PT-C405-pneu, PT-C405-furar, PT-C405-furo, PT-C405-jante, PT-C405-triangulo

\chapter{Um pneu furado}
\label{ch:pt-um-pneu-furado}
\hlchaptermodality{405}

\begin{chapteropening}
\textbf{By the end of this chapter:} \emph{I can deal with a flat tyre: the tyre, getting a puncture, the hole, the rim and the warning triangle.}

Complete the last lesson of chapter 405: I can deal with a flat tyre: the tyre, getting a puncture, the hole, the rim and the warning triangle.
\end{chapteropening}

\section[o pneu]{o pneu — the tyre}

\label{lesson:PT-C405-pneu}



\begin{quote}
Before the new one: say the Portuguese for to fill up, then the Portuguese for to fill to the top.
\end{quote}

\subsection*{You'll want to know: o pneu}
\textbf{o pneu} — \textquotedblleft{}the tyre\textquotedblright{}.

\begin{quote}
O pneu da frente está vazio.
\end{quote}

\begin{quote}
\emph{The front tyre is flat.}
\end{quote}

\emph{Um pneu furado} — a flat tyre, a puncture.

\begin{cousinweb}
\emph{pneu} is short for \emph{pneumático}, a learned borrowing from Latin \emph{pneumaticus}, from Greek \emph{pneumatikós} 'of wind or air' (EtymDB): a tyre is full of air.
\end{cousinweb}

\begin{grammarlens}[title={What you've built}]
One more word for telling, explaining and describing at B1.
\end{grammarlens}

\subsection*{Your turn}
\begin{itemize}
\raggedright
  \item \emph{Say it:} \emph{o pneu}
  \item \emph{Say it:} \emph{o pneu}, once more
  \item \emph{Say it:} a sentence of your own with \emph{o pneu}
  \item \emph{Recall:} say the Portuguese for to fill up, then the Portuguese for to fill to the top, then say \emph{o pneu} again
\end{itemize}

\begin{tcolorbox}[breakable,colback=teal!4,colframe=teal!35!black,title={Before you move on}]
What does o pneu mean? (\textquotedblleft{}The tyre\textquotedblright{}.) Say it once more.
\end{tcolorbox}
\section[furar]{furar — to puncture, to make a hole in}

\label{lesson:PT-C405-furar}



\begin{quote}
Before the new one: say the Portuguese for to fill to the top, then the Portuguese for the tyre.
\end{quote}

\subsection*{You'll want to know: furar}
\textbf{furar} — \textquotedblleft{}to puncture, to make a hole in\textquotedblright{}.

\begin{quote}
Um prego furou o pneu.
\end{quote}

\begin{quote}
\emph{A nail punctured the tyre.}
\end{quote}

\emph{Furar a fila} — to jump the queue (informal).

\begin{cousinweb}
\emph{furar} is inherited from Latin \emph{forāre} 'to bore, to pierce' (EtymDB; Lewis \& Short).
\end{cousinweb}

\begin{grammarlens}[title={What you've built}]
One more word for telling, explaining and describing at B1.
\end{grammarlens}

\subsection*{Your turn}
\begin{itemize}
\raggedright
  \item \emph{Say it:} \emph{furar}
  \item \emph{Say it:} \emph{furar}, once more
  \item \emph{Say it:} a sentence of your own with \emph{furar}
  \item \emph{Recall:} say the Portuguese for to fill to the top, then the Portuguese for the tyre, then say \emph{furar} again
\end{itemize}

\begin{tcolorbox}[breakable,colback=teal!4,colframe=teal!35!black,title={Before you move on}]
What does furar mean? (\textquotedblleft{}To puncture, to make a hole in\textquotedblright{}.) Say it once more.
\end{tcolorbox}
\section[o furo]{o furo — the puncture; the hole}

\label{lesson:PT-C405-furo}



\begin{quote}
Before the new one: say the Portuguese for the tyre, then the Portuguese for to puncture.
\end{quote}

\subsection*{You'll want to know: o furo}
\textbf{o furo} — \textquotedblleft{}the puncture; the hole\textquotedblright{}.

\begin{quote}
Tivemos um furo na estrada para a praia.
\end{quote}

\begin{quote}
\emph{We had a puncture on the road to the beach.}
\end{quote}

In journalism, \emph{um furo} is a scoop.

\begin{cousinweb}
\emph{furo} is built on \emph{furar} 'to pierce', which is inherited from Latin \emph{forāre} 'to bore' (Lewis \& Short).
\end{cousinweb}

\begin{grammarlens}[title={What you've built}]
One more word for telling, explaining and describing at B1.
\end{grammarlens}

\subsection*{Your turn}
\begin{itemize}
\raggedright
  \item \emph{Say it:} \emph{o furo}
  \item \emph{Say it:} \emph{o furo}, once more
  \item \emph{Say it:} a sentence of your own with \emph{o furo}
  \item \emph{Recall:} say the Portuguese for the tyre, then the Portuguese for to puncture, then say \emph{o furo} again
\end{itemize}

\begin{tcolorbox}[breakable,colback=teal!4,colframe=teal!35!black,title={Before you move on}]
What does o furo mean? (\textquotedblleft{}The puncture; the hole\textquotedblright{}.) Say it once more.
\end{tcolorbox}
\section[a jante]{a jante — the wheel rim}

\label{lesson:PT-C405-jante}



\begin{quote}
Before the new one: say the Portuguese for to puncture, then the Portuguese for the puncture.
\end{quote}

\subsection*{You'll want to know: a jante}
\textbf{a jante} — \textquotedblleft{}the wheel rim\textquotedblright{}.

\begin{quote}
A jante ficou riscada quando bati no passeio.
\end{quote}

\begin{quote}
\emph{The rim got scratched when I hit the kerb.}
\end{quote}

In Brazil: \emph{o aro} or \emph{a roda}.

\begin{cousinweb}
\emph{jante} was borrowed from French \emph{jante} 'wheel rim', which EtymDB traces to Vulgar Latin \emph{gambita}; the earlier history of that word is not settled.
\end{cousinweb}

\begin{grammarlens}[title={What you've built}]
One more word for telling, explaining and describing at B1.
\end{grammarlens}

\subsection*{Your turn}
\begin{itemize}
\raggedright
  \item \emph{Say it:} \emph{a jante}
  \item \emph{Say it:} \emph{a jante}, once more
  \item \emph{Say it:} a sentence of your own with \emph{a jante}
  \item \emph{Recall:} say the Portuguese for to puncture, then the Portuguese for the puncture, then say \emph{a jante} again
\end{itemize}

\begin{tcolorbox}[breakable,colback=teal!4,colframe=teal!35!black,title={Before you move on}]
What does a jante mean? (\textquotedblleft{}The wheel rim\textquotedblright{}.) Say it once more.
\end{tcolorbox}
\section[o triângulo]{o triângulo — the warning triangle; the triangle}

\label{lesson:PT-C405-triangulo}



\begin{quote}
Before the new one: say the Portuguese for the puncture, then the Portuguese for the wheel rim.
\end{quote}

\subsection*{You'll want to know: o triângulo}
\textbf{o triângulo} — \textquotedblleft{}the warning triangle; the triangle\textquotedblright{}.

\begin{quote}
Ele pôs o triângulo atrás do carro parado.
\end{quote}

\begin{quote}
\emph{He put the warning triangle behind the stopped car.}
\end{quote}

In Portugal a car must carry \emph{o triângulo} and a reflective vest.

\begin{cousinweb}
\emph{triângulo} is a learned borrowing from Latin \emph{triangulus} 'three-cornered', from \emph{trēs} 'three' and \emph{angulus} 'corner' (EtymDB).
\end{cousinweb}

\begin{grammarlens}[title={What you've built}]
One more word for telling, explaining and describing at B1.
\end{grammarlens}

\subsection*{Your turn}
\begin{itemize}
\raggedright
  \item \emph{Say it:} \emph{o triângulo}
  \item \emph{Say it:} \emph{o triângulo}, once more
  \item \emph{Say it:} a sentence of your own with \emph{o triângulo}
  \item \emph{Recall:} say the Portuguese for the puncture, then the Portuguese for the wheel rim, then say \emph{o triângulo} again
\end{itemize}

\begin{tcolorbox}[breakable,colback=teal!4,colframe=teal!35!black,title={Before you move on}]
What does o triângulo mean? (\textquotedblleft{}The warning triangle; the triangle\textquotedblright{}.) Say it once more.
\end{tcolorbox}
